% !TEX root = ../main.tex
\chapter{État de l'art et Revue de la Littérature}
\label{chap:etat_art}

\section{Généralités sur la Commande Numérique par Calculateur (CNC)}
La Commande Numérique par Calculateur (CNC) se définit comme le pilotage automatisé des mouvements et des fonctions auxiliaires d'une machine-outil par microprocesseur, via l'interprétation d'un programme d'instructions textuelles (G-Code). Depuis les travaux pionniers de John T. Parsons (qui a utilisé des cartes perforées pour usiner des pales d'hélicoptère dans les années 1940), la CNC a évolué de systèmes câblés immenses et rigides vers des architectures logicielles distribuées, flexibles et miniaturisées.

\subsection{Évolution historique des systèmes CNC}
L'histoire de la commande numérique est intrinsèquement liée à l'évolution de l'informatique industrielle. Le tableau~\ref{tab:histoire_cnc} retrace les grandes ruptures technologiques qui ont conduit aux systèmes modernes ouverts.

\begin{table}[htbp]
\centering
\renewcommand{\arraystretch}{1.3}
\begin{tabular}{|c|p{2.5cm}|p{6cm}|p{4.5cm}|}
\hline
\rowcolor{macouleur!20}
\textbf{Décennie} & \textbf{Technologie} & \textbf{Innovation principale} & \textbf{Limites identifiées} \\
\hline
1940--1950 & Cartes perforées & Travaux de Parsons / MIT : premier usinage de pale d'hélicoptère & Pas de rétroaction, séquentiel pur \\
\hline
1960--1970 & NC câblé & Contrôleurs à relais, tubes à vide & Rigides, non-reprogrammables \\
\hline
1970--1980 & CNC 8 bits & Microprocesseur Intel 8080 / Z80, DNC & Un seul cœur, pas de multi-axes \\
\hline
1980--1990 & Propriétaire & FANUC, Siemens, Heidenhain — OS fermés & Coût prohibitif, dépendance \\
\hline
1990--2010 & PC-Based (Linux) & Interface PC + carte FPGA (Mesa Cards) & Configuration complexe, coût \\
\hline
2015--2024 & RTOS 32-bits open & ESP32, STM32 + FreeRTOS, GRBL-ESP32 & Notre domaine d'innovation \\
\hline
\end{tabular}
\caption{Évolution historique des architectures de commande numérique}
\label{tab:histoire_cnc}
\end{table}

\subsection{Principe d'asservissement d'un système CNC}
Un système CNC moderne s'articule autour d'une chaîne de commande comportant trois unités principales agissant en synergie :
\begin{enumerate}
    \item \textbf{L'Unité de Traitement des Données (Data Processing Unit -- DPU) :} Elle correspond au logiciel IHM. Elle lit le fichier de commande numérique, l'analyse, simule la trajectoire et l'envoie de manière séquencée.
    \item \textbf{L'Unité de Contrôle (Control Loop Unit -- CLU) :} C'est le « cerveau interne » (le microcontrôleur sur la machine). Il reçoit les données du DPU, calcule la cinématique, gère les accélérations (planner) et génère les signaux électriques logiques (Step/Direction ou signaux PWM).
    \item \textbf{L'Unité de Puissance et d'Actionnement :} Composée de drivers (variateurs) qui amplifient les signaux logiques de la CLU, et de moteurs électriques transmettant les couples et mouvements aux axes physiques.
\end{enumerate}

\section{Classification Architecturale des Machines CNC 5 Axes}
Le passage d'une machine 3 axes à une machine 5 axes introduit une complexité majeure dans la chaîne cinématique. Pour qu'une fraise puisse atteindre n'importe quel point d'une pièce sous n'importe quel angle, deux axes de rotation (généralement nommés A, B ou C) doivent s'ajouter aux trois translations (X, Y, Z). Il existe trois grandes familles d'architectures pour les machines 5 axes, classées selon l'emplacement de ces deux axes rotatifs.

\begin{figure}[htbp]
    \centering
    \includegraphics[width=0.6\textwidth]{images/trunnion_diagram.png}
    \caption{Schéma cinématique d'une machine CNC 5 axes de type Trunnion (Table-Table)}
    \label{fig:trunnion_schema}
\end{figure}

\begin{figure}[htbp]
    \centering
    \begin{tikzpicture}[
        node distance=2.5cm and 1.8cm,
        box/.style={rectangle, draw, rounded corners=4pt, align=center, fill=blue!5, minimum width=3.8cm, minimum height=1.8cm, thick},
        title/.style={font=\bfseries\small},
        desc/.style={font=\footnotesize, text width=3.8cm, align=center},
        arrow/.style={-{Stealth[scale=1.2]}, thick, color=gray!80}
    ]
        \node[box] (table_table) {
            \textbf{Table-Table}\\(Trunnion)
        };
        \node[below=0.2cm of table_table, desc, align=center] {Rotations A et C sur la table.\\ Meilleure rigidité.};
        
        \node[box, right=1cm of table_table] (head_head) {
            \textbf{Head-Head}\\(Tête Pivotante)
        };
        \node[below=0.2cm of head_head, desc, align=center] {Rotations A et C sur la broche.\\ Pièces volumineuses.};
        
        \node[box, right=1.5cm of head_head] (head_table) {
            \textbf{Head-Table}\\(Hybride)
        };
        \node[below=0.2cm of head_table, desc] {Répartition des rotations\\(mixte Tête / Table).};
        
        \draw[arrow] (table_table.east) -- (head_head.west) node[midway, above, font=\scriptsize, color=black] {Cinématique};
        \draw[arrow] (head_head.east) -- (head_table.west) node[midway, above, font=\scriptsize, color=black] {Compromis};
    \end{tikzpicture}
    \caption{Comparaison schématique des topologies CNC 5 axes}
    \label{fig:topo_5axes}
\end{figure}

\subsection{Configuration «Table-Table» (Trunnion Table)}
Dans cette architecture, les deux axes rotatifs (souvent l'axe A pour le basculement et l'axe B pour la rotation continue) sont emboîtés l'un dans l'autre et supportent la table de travail.
\begin{itemize}
    \item \textbf{Cinématique :} L'outil se déplace linéairement en X, Y et Z. C'est la pièce à usiner qui tourne (Pan) et bascule (Tilt) face à l'outil.
    \item \textbf{Avantages :} Meilleure rigidité globale et excellente précision sur les petites pièces compactes.
    \item \textbf{Limites :} Le volume de la pièce et son poids sont strictement limités par le diamètre du plateau rotatif.
\end{itemize}
C'est l'architecture que nous avons retenue pour notre projet didactique.

\subsection{Configuration «Head-Head» (Tête Pivotante Bi-Rotative)}
Ici, la table supportant la pièce est fixe. Les deux axes rotatifs sont intégrés dans la tête qui supporte la broche.
\begin{itemize}
    \item \textbf{Avantages :} Parfaitement adaptée à l'usinage de pièces de très grandes dimensions et très lourdes.
    \item \textbf{Limites :} Rigidité angulaire diminuée. Le bras de levier important amplifie considérablement la moindre erreur angulaire.
\end{itemize}

\subsection{Configuration «Head-Table» (Hybride)}
Cette architecture répartit les rotations. Généralement, l'axe de basculement est fixé sur la tête de la broche, et l'axe de rotation continue est placé sur une table tournante. Cette solution hybride permet un compromis intéressant pour les machines de taille moyenne.

\section{Comparatif des Machines CNC 5 Axes Accessibles}
Afin de situer notre projet dans le paysage industriel actuel, le tableau~\ref{tab:comparatif_cnc} compare les caractéristiques de notre prototype avec celles de machines de référence.

\begin{table}[htbp]
\centering
\renewcommand{\arraystretch}{1.3}
\begin{tabular}{|p{2.5cm}|p{2cm}|p{2cm}|p{2.5cm}|p{2.5cm}|}
\hline
\rowcolor{macouleur!20}
\textbf{Critère} & \textbf{PFE (Nous)} & \textbf{PocketNC} & \textbf{Haas UMC-500} & \textbf{Candle+GRBL} \\
\hline
\textbf{Architecture} & Trunnion A+B & Trunnion A+B & Trunnion A+C & 3 axes \\
\hline
\textbf{Volume} & 20cm cube & 12cm cube & 50cm cube & Variable \\
\hline
\textbf{Contrôleur} & ESP32 RTOS & Propriétaire & Haas CNC & Arduino \\
\hline
\textbf{Précision} & $\pm 0{,}05$ mm & $\pm 0{,}025$ mm & $\pm 0{,}005$ mm & $\pm 0{,}1$ mm \\
\hline
\textbf{Open-Source} & Oui (100\%) & Non & Non & Oui (firm) \\
\hline
\textbf{Coût} & $\approx 0{,}5 M$ & $\approx 7 M$ & $\approx 300 M$ & $\approx 0{,}1 M$ \\
\hline
\textbf{IHM} & Flutter App & Kinetic & Tactile HAAS & PC Candle \\
\hline
\end{tabular}
\caption{Tableau comparatif des machines CNC 5 axes du marché vs. notre prototype}
\label{tab:comparatif_cnc}
\end{table}

L'analyse de ce tableau démontre que notre prototype, bien qu'inférieur en précision et en volume aux machines industrielles, se positionne comme la \textbf{seule solution entièrement open-source} dotée d'une IHM multi-plateforme moderne (Flutter) et d'une connectivité sans fil. Le rapport \textbf{performance/coût} est sans équivalent dans le segment éducatif africain.

\section{Problématiques Avancées : Singularités et TCP}
L'usinage 5 axes introduit des défis théoriques absents des machines cartésiennes classiques : la gestion des singularités et le contrôle du point de contact outil.

\subsection{Les Singularités Cinématiques (Gimbal Lock)}
Une singularité cinématique survient lorsque deux axes de rotation s'alignent, entraînant la perte d'un degré de liberté. Pour une table Trunnion (Axe A et B), cela se produit généralement lorsque l'Axe B est vertical et que l'angle de l'Axe A atteint une valeur critique (souvent $0^\circ$).
À ce point précis, une infinité de combinaisons d'angles peut donner la même position spatiale, ce qui peut provoquer des accélérations infinies théoriques des moteurs et des marques de vibration sur la pièce (\textit{chatter}). Le logiciel de contrôle doit intégrer des algorithmes d'évitement ou de lissage de trajectoire aux abords de ces zones critiques.

Il existe trois types principaux de singularités rencontrées en usinage 5 axes :
\begin{enumerate}
    \item \textbf{Singularité de configuration (Gimbal Lock) :} Perte d'un degré de liberté par alignement des axes de rotation. Se manifeste par des vitesses angulaires théoriquement infinies.
    \item \textbf{Singularité de frontière :} Le vecteur outil atteint la limite physique d'une articulation (débattement angulaire maximal). Un interverrouillage logiciel est nécessaire.
    \item \textbf{Singularité de chemin :} Le parcours CAO génère une trajectoire passant exactement par un point singulier. Le post-processeur FAO doit détecter ces cas et déranger légèrement la trajectoire.
\end{enumerate}

\subsection{Le Concept de TCP (Tool Center Point Control)}
Dans un système 3 axes, la position de l'outil est directement liée aux coordonnées X, Y, Z. En 5 axes, toute rotation de la table (Axe A ou B) déplace physiquement la pièce par rapport à la pointe de l'outil.
Le \textbf{TCP Control} est la capacité du contrôleur à compenser automatiquement ces décalages géométriques. Au lieu de programmer les positions des moteurs, le programmeur définit la trajectoire de la \textbf{pointe de l'outil} par rapport à la pièce. C'est l'intelligence embarquée (le microcontrôleur) qui recalcule en temps réel les compensations linéaires $\Delta X, \Delta Y, \Delta Z$ nécessaires pour maintenir l'outil sur le point de contact nominal avec la surface.

\section{L'Écosystème Logiciel et le Langage G-Code 5 Axes}

\subsection{La Chaîne Numérique 5 Axes}
Le flux de données d'une pièce usinée numériquement passe par trois phases :
\begin{itemize}
    \item \textbf{CAO (Conception Assistée par Ordinateur) :} Le concepteur crée le modèle 3D de la pièce.
    \item \textbf{FAO (Fabrication Assistée par Ordinateur) :} Le programmeur génère les parcours d'outil. En 5 axes, on distingue deux méthodes : le \textit{3+2 axes} (Usinage Positionnel) et le \textit{5 axes continus} (simultané).
    \item \textbf{Le Post-Processeur :} Il traduit la géométrie (vecteurs de l'outil I, J, K) en positions angulaires réelles (A, B) des moteurs de la machine spécifique.
\end{itemize}

\subsection{Le Langage G-Code (Norme ISO 6983)}
C'est le langage universel de la commande numérique. Un programme est une suite alphabétique de blocs d'instructions. En multi-axes, un bloc ressemble à ceci :
\texttt{G01 X15.5 Y0.0 Z-5.0 A45.0 B90.0 F800}

Le tableau~\ref{tab:gcode_5axes} présente les codes G et M principaux utilisés dans notre logiciel :

\begin{table}[htbp]
\centering
\renewcommand{\arraystretch}{1.3}
\begin{tabular}{|l|p{6cm}|l|}
\hline
\rowcolor{macouleur!20}
\textbf{Code} & \textbf{Fonction} & \textbf{Utilisation} \\
\hline
G00 & Déplacement rapide (Rapid) en XYZ & Positionnement \\
\hline
G01 & Interpolation linéaire à vitesse contrôlée & Usinage rectiligne \\
\hline
G02 & Interpolation circulaire sens horaire & Congé, rayon \\
\hline
G03 & Interpolation circulaire sens anti-horaire & Congé, rayon \\
\hline
G17/18/19 & Sélection du plan de travail (XY, XZ, YZ) & Fraisage de poches \\
\hline
G28 & Retour à l'origine machine (Homing) & Initialisation \\
\hline
G54--G59 & Systèmes de coordonnées pièce & Multi-fixtures \\
\hline
G91 & Mode de programmation incrémentale & Jog relatif \\
\hline
G92 & Définition de l'origine pièce courante & Palpage TCP \\
\hline
M03 / M04 & Mise en marche de la broche (sens direct / inverse) & Fraisage \\
\hline
M05 & Arrêt de la broche & Fin pass \\
\hline
M08 / M09 & Arrosage ON / OFF & Refroidissement \\
\hline
M30 & Fin de programme + retour au début & Fin cycle \\
\hline
\$H & Homing ESP32 (FluidNC/GRBL étendu) & Calibration \\
\hline
\$X & Déverrouillage de l'état ALARM & Reprise \\
\hline
\end{tabular}
\caption{Tableau des principaux codes G et M supportés par notre contrôleur ESP32}
\label{tab:gcode_5axes}
\end{table}

\section{Analyse Comparative des Protocoles de Communication Industrielle}

Le choix du protocole de communication entre l'IHM Flutter et le microcontrôleur ESP32 est central pour les performances temps réel du système. Le tableau~\ref{tab:protocoles} compare les principales solutions envisageables.

\begin{table}[htbp]
\centering
\renewcommand{\arraystretch}{1.3}
\resizebox{\textwidth}{!}{
\begin{tabular}{|l|c|c|c|c|p{4cm}|}
\hline
\rowcolor{macouleur!20}
\textbf{Protocole} & \textbf{Latence} & \textbf{Débit max} & \textbf{Connectivité} & \textbf{Overhead} & \textbf{Adéquation CNC} \\
\hline
\textbf{USB-Serial (GRBL)} & $< 10$ ms & 115 200 baud & Filaire & Faible & Limité (câble obligatoire) \\
\hline
\textbf{TCP (WebSocket)} & $5$--$50$ ms & 100 Mbps & Wi-Fi/Eth & Moyen & Bon (handshake) \\
\hline
\textbf{UDP (choix retenu)} & $1$--$5$ ms & 100 Mbps & Wi-Fi & Très faible & \textbf{Optimal} (pas d'ACK) \\
\hline
\textbf{Modbus RTU} & $10$--$100$ ms & 1 Mbps & RS-485 & Moyen & Industriel, trop lent \\
\hline
\textbf{EtherCAT} & $< 0{,}1$ ms & 100 Mbps & Ethernet & Faible & Professionnel, complexe \\
\hline
\textbf{Profibus} & $1$--$10$ ms & 12 Mbps & RS-485 & Élevé & Obsolète, propriétaire \\
\hline
\end{tabular}
}
\caption{Tableau comparatif des protocoles de communication pour systèmes CNC}
\label{tab:protocoles}
\end{table}

\textbf{Justification du choix UDP :} La bande passante théorique nécessaire pour notre application est calculée comme suit. En supposant un taux de rafraîchissement de $f_{refresh} = 50$ trames/seconde et une taille de trame JSON compressée de $T_{trame} = 128$ octets :
\begin{equation}
    BW_{required} = N_{trames} \times T_{trame} \times 8 = 50 \times 128 \times 8 = \boxed{51~200 \text{ bps} \approx 51 \text{ kbps}}
\end{equation}
Ce débit est négligeable face aux 100 Mbps disponibles en Wi-Fi 802.11n. L'UDP est donc choisi pour sa latence minimale ($< 2$ ms) et l'absence d'overhead d'accusé de réception (ACK), critique lors de l'envoi en rafale des blocs G-Code.

\section{État de l'Art des Firmwares CNC Open-Source}

L'intelligence de la machine réside dans son micrologiciel (firmware). Pour opérer une machine 5 axes open-source, plusieurs grands firmwares existent. Le tableau~\ref{tab:firmwares} les compare.

\begin{table}[htbp]
\centering
\renewcommand{\arraystretch}{1.3}
\resizebox{\textwidth}{!}{
\begin{tabular}{|l|c|c|c|p{5cm}|}
\hline
\rowcolor{macouleur!20}
\textbf{Firmware} & \textbf{Contrôleur} & \textbf{Axes max} & \textbf{5 axes continus} & \textbf{Limitation principale} \\
\hline
\textbf{GRBL (classique)} & Arduino Uno (AVR 8-bit) & 3 & Non & Insufficient Timers. Inadapté. \\
\hline
\textbf{Marlin 2.x} & AVR 32-bit / ARM & 6 & Limité & Conçu pour impression 3D, pas pour fraisage. \\
\hline
\textbf{Klipper} & Raspberry Pi + MCU & 6 & Partiel & Complexe, nécessite Linux et Python. \\
\hline
\textbf{FluidNC / GRBL-ESP32} & ESP32 & 6 & Oui (avec plugin) & Configuration YAML ; pas de cinématique Trunnion native. \\
\hline
\textbf{LinuxCNC} & PC + FPGA Mesa & Illimité & Oui & Très complexe, onéreux, nécessite un PC dédié sous Linux RT. \\
\hline
\textbf{Notre firmware (Custom)} & ESP32 Dual-Core & 5 & Oui & Spécifique à l'architecture Trunnion de ce PFE (avantage). \\
\hline
\end{tabular}
}
\caption{Comparatif des firmwares CNC open-source pour machine 5 axes}
\label{tab:firmwares}
\end{table}

Ce tableau confirme la nécessité d'un développement ``from scratch'' du firmware ESP32. Aucun des firmwares existants ne gère nativement la cinématique inverse d'une table Trunnion (axes A+B emboîtés) avec compensation TCP en temps réel.

\section{Matériaux Usinables et Usinabilité}

Le choix des matériaux d'essai est guidé par leur usinabilité et leur pertinence industrielle dans le contexte malien. Le tableau~\ref{tab:usinabilite} présente les paramètres de coupe recommandés pour notre broche de 500\,W et notre fraise de $6\,\text{mm}$.

\begin{table}[htbp]
\centering
\renewcommand{\arraystretch}{1.3}
\resizebox{\textwidth}{!}{
\begin{tabular}{|l|c|c|c|c|c|c|}
\hline
\rowcolor{macouleur!20}
\textbf{Matériau} & \textbf{Dureté (HB)} & \textbf{$k_c$ (N/mm²)} & \textbf{$V_c$ (m/min)} & \textbf{$f_z$ (mm/dent)} & \textbf{$a_p$ max (mm)} & \textbf{Indice Usinab.} \\
\hline
Résine Epoxy (Cibatool) & 15--25 & 200 & 300--500 & 0,10 & 3,0 & Très Facile \\
\hline
Bois dur (Chêne, Iroko) & 25--35 & 350 & 200--400 & 0,12 & 4,0 & Facile \\
\hline
HDPE (Polyéthylène HD) & 60--70 & 400 & 200--350 & 0,08 & 3,0 & Facile \\
\hline
Aluminium AW-2017A & 95--105 & 700 & 100--200 & 0,05 & 1,0 & Moyen \\
\hline
Aluminium AW-5052 & 60 & 600 & 150--250 & 0,06 & 1,5 & Moyen \\
\hline
Laiton CuZn37 & 70--80 & 800 & 80--150 & 0,04 & 0,8 & Moyen \\
\hline
Cuivre pur & 45--55 & 600 & 80--120 & 0,05 & 0,8 & Difficile (adhérent) \\
\hline
Acier doux (S235) & 120 & 1\,500 & 20--50 & 0,02 & 0,3 & Hors capacité broche \\
\hline
\end{tabular}
}
\caption{Paramètres d'usinabilité et vitesses de coupe pour les matériaux testables avec notre CNC}
\label{tab:usinabilite}
\end{table}

\section{Conclusion du chapitre}

Ce chapitre a établi le socle théorique et comparatif sur lequel repose notre démarche de conception. Le passage de 3 à 5 axes exige une refonte totale de la géométrie de la machine (table Trunnion), des solutions de transmissions strictes (Vis à billes), et surtout, d'une rupture architecturale en matière de puissance de calcul embarquée (ESP32 Dual-Core). L'analyse comparative des protocoles de communication, des firmwares open-source et des matériaux usinables a permis de valider et justifier chaque choix technologique majeur du projet. Le chapitre suivant s'attachera à formaliser précisément ces besoins et à déterminer les calculs de sollicitations mécaniques de la future CNC 5 axes.
